\section{FAQ y síntesis de la experiencia}

\subsection{Preguntas frecuentes}

\begin{itemize}
    \item \textbf{P: ¿Para qué tamaño de proyecto es adecuado DSPy?} R: Es adecuado para todos los proyectos que necesitan combinar múltiples componentes de LLM, sobre todo cuando se quiere gobernar el prompt mediante código.
    \item \textbf{P: ¿Durante cuánto tiempo hay que conservar la Evidence Matrix?} R: Al menos 30 días o la ventana de tiempo del rollout anterior, para poder hacer regress reviews.
    \item \textbf{P: ¿Cómo se maneja el optimizer drift?} R: Primero se revisa la Evidence Matrix de manera retrospectiva, luego se ejecuta una prueba comparativa con el fallback prompt, y al mismo tiempo se registra el metric delta.
    \item \textbf{P: ¿Se necesita un DSL especializado?} R: DSPy ya incluye su propio DSL: basta con declarar los módulos y la signature en Python, sin necesidad de reinventar la rueda.
    \item \textbf{P: ¿Cómo se mide el optimizer improvements?} R: Se cuantifica mediante la combinación de tres indicadores principales: `automation coverage`, `latency` y `metric accuracy`.
\end{itemize}

\subsection{Lecciones extraídas de la experiencia}

La experiencia nos dice: 1) en cada release hay que entregar el evidence snapshot a QA; 2) el monitor necesita 24/7 alerting + weekly review; 3) la documentation no puede quedarse solo en la cabeza, debe escribirse como `DSPy.md` + `Incident Log`.

\subsection{Resumen de la sección}

El FAQ y la síntesis de la experiencia le dan al equipo respuestas rápidas a las preguntas frecuentes sobre el pipeline de DSPy, y evitan que se repitan los mismos errores en la práctica.

